\clearpage\thispagestyle{empty}
\pdfbookmark[0]{Quellen zur Bestimmung des Osterfestes (1807)}{osterfest1807-sources}
\begin{center}\textbf{Quellen zur Bestimmung des Osterfestes (1807)}\end{center}
\noindent Diese KI-Übertragung des deutschen Originals in heutiges Deutsch folgt Gauß’ \textit{Noch etwas über die Bestimmung des Osterfestes} in den \textit{Werken}, Band~VI (Göttingen, 1874), S.~82--86. Der Erstdruck erschien im \textit{Braunschweigischen Magazin}, Band~20, 12.~September 1807, Sp.~589--596, unter dem Namen Doktor Gauß. \href{https://www.digitale-sammlungen.de/de/view/bsb10531610?page=303}{Erstdruck}. \href{https://gdz.sub.uni-goettingen.de/id/PPN236007467}{Historischer deutscher Text in den Werken}.
\par\smallskip\noindent Im Rechenbeispiel für 1801 stehen in der Tabelle zur vierten Regel im Erstdruck 90 und 7300 (Sp.~592), in den \textit{Werken} 48 und 7258 (S.~84). Beide Drucke geben 6 als den daraus folgenden vierten Rest an. Die Verweise in eckigen Klammern auf S.~75, 78 und 80 stammen aus den \textit{Werken}. Zu den verschiedenen Fassungen der Osterregeln siehe Reinhold Bien, \href{https://doi.org/10.1007/s00407-004-0078-5}{\textit{Gauß and Beyond: The Making of Easter Algorithms}}, \textit{Archive for History of Exact Sciences} 58 (2004), S.~439--452, besonders S.~441--443.
\par\smallskip\noindent Übertragung und Satz wurden mit KI erstellt.
\clearpage
